% Mechanically converted from the audited Markdown manuscript.
\documentclass[11pt]{article}
\usepackage[a4paper,margin=27mm]{geometry}
\usepackage[T1]{fontenc}
\usepackage[utf8]{inputenc}
\usepackage{lmodern}
\usepackage{microtype}
\usepackage{amsmath,amssymb,amsthm,mathtools}
\usepackage{booktabs}
\usepackage{xcolor}
\usepackage[colorlinks=true,linkcolor=blue!55!black,citecolor=blue!55!black,urlcolor=blue!65!black]{hyperref}
\usepackage{fancyhdr}

\numberwithin{equation}{section}
\pagestyle{fancy}
\fancyhf{}
\fancyhead[L]{Carptopus}
\fancyhead[R]{Higher-order Stirling cycle rows}
\fancyfoot[C]{\thepage}
\setlength{\headheight}{14pt}
\setlength{\parindent}{0pt}
\setlength{\parskip}{0.42em}
\setlength{\emergencystretch}{4em}

\title{A complete log-concavity classification\\
for higher-order Stirling cycle rows}
\author{Carptopus\\\href{mailto:carptopus@163.com}{carptopus@163.com}}
\date{Manuscript, 2 September 2026}

\hypersetup{
  pdftitle={A complete log-concavity classification for higher-order Stirling cycle rows},
  pdfauthor={Carptopus},
  pdfsubject={Log-concavity of higher-order associated Stirling cycle rows},
  pdfkeywords={associated Stirling numbers, log-concavity, recurrence, Fourier--Edgeworth estimate, computer-assisted proof}
}

\begin{document}
\maketitle

\begin{abstract}

Let $[N,k]_{\ge r}$ denote the number of permutations of $N$ elements having exactly $k$
cycles, every cycle of length at least $r$, and define the higher-order Stirling cycle rows

\[
C_r(n,k)=[n+(r-1)k,k]_{\ge r},\qquad 0\le k\le n.
\]

Deb and Sokal proved that these rows are log-concave for $r=2$ and asked whether the same
holds for $r=3,4,5$. We prove all three remaining cases. Together with the elementary
$r=1$ case, this gives a complete classification: every row $C_r(n,\cdot)$ is log-concave
for every $n$ if and only if $1\le r\le5$. The cutoff is already visible in the row $n=3$;
an explicit ratio is strictly decreasing in $r$, is greater than one at $r=5$, and is less
than one at $r=6$.

The proofs use three different mechanisms. The case $r=3$ follows from Sagan's inductive
criterion. For $r=4$ we introduce a factorial normalization and prove that the recurrence
preserves the narrower weighted log-concavity cone determined by the removed factorials.
For $r=5$ we combine two exact parameter wedges, a certified finite prefix of $5,754,989$
inequalities, and an effective fifth-order Fourier--Edgeworth estimate on the remaining
compact band. We also prove a uniform structural result valid for every $r\ge2$: the
normalized rows

\[
\frac{k!\,C_r(n,k)}{(n+(r-1)k)!}
\]

are log-concave.

\end{abstract}

\section{Introduction}


For integers $r\ge1$, let $[N,k]_{\ge r}$ be the signless $r$-associated Stirling number
of the first kind. Thus $[N,k]_{\ge r}$ counts permutations of an $N$-element set having
$k$ cycles, with every cycle of length at least $r$. The higher-order Stirling cycle
triangle studied by Deb and Sokal \cite{deb-sokal} has entries

\[
C_r(n,k)=[n+(r-1)k,k]_{\ge r}.
\tag{1.1}
\]

The unusual feature of (1.1) is that the size of the ground set changes with $k$ along a
fixed row. Consequently, real-rootedness results for the polynomial
$\sum_k[N,k]_{\ge r}x^k$ with fixed $N$ do not imply log-concavity of the rows in (1.1).
Indeed, Deb and Sokal showed that the row-generating polynomials in (1.1) generally have
nonreal zeros when $r\ge3$ \cite{deb-sokal}.

Deb and Sokal proved row log-concavity for $r=2$ and conjectured it for $r=3,4,5$
\cite[Conjecture 1.3(c)]{deb-sokal}. The same three cases were recorded as Problem 30.9 in the BCC30
problem collection \cite{bcc30}. Our main result settles those remaining parameter layers and also
locates the exact cutoff.

\textbf{Theorem 1.1 (complete classification).} For an integer $r\ge1$, the sequence

\[
C_r(n,0),C_r(n,1),\ldots,C_r(n,n)
\]

is log-concave for every $n\ge0$ if and only if $1\le r\le5$.

Here and below, entries outside $0\le k\le n$ are zero. The positive part of each row is
contiguous, so the usual adjacent inequalities are sufficient.

The proof naturally separates into four pieces. The ordinary case $r=1$ follows from

\[
\sum_k [n,k]_{\ge1}x^k=x(x+1)\cdots(x+n-1),
\]

and $r=2$ is due to Deb and Sokal. Section 3 proves $r=3$ by applying Sagan's
log-concavity-preserving recurrence criterion \cite{sagan}. Section 4 proves $r=4$ by preserving a
weighted cone rather than the full cone of log-concave sequences. Sections 6--10 give the
$r=5$ proof. Section 5 shows that every $r\ge6$ fails already for $n=3$.

One structural result applies at every order and explains why factorial normalization is
useful even when the original rows cease to be log-concave.

\textbf{Theorem 1.2 (normalized rows).} Let $r\ge2$ and put

\[
H_r(n,k)=\frac{k!\,C_r(n,k)}{(n+(r-1)k)!}.
\tag{1.2}
\]

For every $n\ge0$, the sequence $H_r(n,0),\ldots,H_r(n,n)$ is log-concave.

The $r=5$ part of Theorem 1.1 is computer-assisted. Its finite and symbolic components use
exact rational arithmetic or decimal arithmetic with a proved relative-rounding bound. The analytic tail is
reduced to explicit rational inequalities. Section 11 records the verification boundary and
the reproducibility entry points.

\section{Recurrences and an all-order normalization}


Write $p=r-1$ and

\[
N=n+pk.
\]

The standard insertion recurrence for associated Stirling numbers becomes

\[
C_r(n,k)
=(N-1)C_r(n-1,k)+(N-1)_{r-1}C_r(n-1,k-1),
\tag{2.1}
\]

where $(x)_j=x(x-1)\cdots(x-j+1)$. The first term inserts the new element into an
existing cycle. The second term forms a new $r$-cycle from the new element and $r-1$
additional elements.

After the normalization (1.2), all falling factorials cancel and (2.1) becomes

\[
H_r(n,k)=a_kH_r(n-1,k)+b_kH_r(n-1,k-1),
\tag{2.2}
\]

with

\[
a_k=\frac{n+pk-1}{n+pk},
\qquad
b_k=\frac{k}{n+pk}.
\tag{2.3}
\]

We use the following form of Sagan's criterion [3, Theorem 1]. Suppose a nonnegative
triangular array satisfies

\[
y_k=d_kx_k+c_kx_{k-1}.
\]

If $c_k$ and $d_k$ are log-concave in $k$ and

\[
2c_kd_k\ge c_{k-1}d_{k+1}+c_{k+1}d_{k-1},
\tag{2.4}
\]

then the transformation preserves log-concavity. Although Sagan states the theorem for
nonnegative integer arrays, its proof uses only addition, multiplication, and nonnegative
inequalities, so it applies without change to nonnegative rational arrays.

\textbf{Proof of Theorem 1.2.} For the coefficients in (2.3), direct simplification gives, for
$n\ge2$ and $1\le k\le n-1$,

\[
a_k^2-a_{k-1}a_{k+1}
=\frac{p^2(2kp+2n-1)}{(kp+n)^2(kp+n-p)(kp+n+p)},
\tag{2.5}
\]

\[
b_k^2-b_{k-1}b_{k+1}
=\frac{n(2kp+n)}{(kp+n)^2(kp+n-p)(kp+n+p)},
\tag{2.6}
\]

and

\[
2a_kb_k-a_{k-1}b_{k+1}-a_{k+1}b_{k-1}
=\frac{2p(knp+kp+n^2)}{(kp+n)^2(kp+n-p)(kp+n+p)}.
\tag{2.7}
\]

All three quantities are nonnegative in the stated range. The zero boundaries are immediate,
and the initial row consists only of $H_r(0,0)=1$. Sagan's criterion applied successively to
(2.2) proves the theorem. $\square$

The removed factor

\[
W_r(n,k)=\frac{(n+pk)!}{k!}
\tag{2.8}
\]

is generally log-convex in $k$. Theorem 1.2 therefore does not imply log-concavity of
$C_r(n,\cdot)$, but it provides the common base sequence used in the $r=4$ argument.

\section{The third-order rows}


\textbf{Theorem 3.1.} For every $n\ge0$, the row $C_3(n,\cdot)$ is log-concave.

\textbf{Proof.} In (2.1), put $r=3$ and

\[
L=n+2k-1.
\]

Then

\[
C_3(n,k)=L C_3(n-1,k)+L(L-1)C_3(n-1,k-1).
\tag{3.1}
\]

In Sagan's notation, take

\[
c_k=L(L-1),\qquad d_k=L.
\]

The sequence $d_k$ is a nonnegative arithmetic progression. The sequence $c_k$ is the
pointwise product of two nonnegative arithmetic progressions. Both are log-concave on the
required index range. The mixed defect in (2.4) is

\[
2d_kc_k-d_{k-1}c_{k+1}-d_{k+1}c_{k-1}=8(L-1)\ge0.
\tag{3.2}
\]

The initial row and the zero boundaries satisfy the claim, so induction using Sagan's
criterion completes the proof. $\square$

\section{The fourth-order rows}


The direct use of Sagan's criterion fails at $r=4$: its mixed coefficient defect has the
wrong sign. We retain the factorial normalization but impose the stronger inequality needed
after restoring the weights.

For this section set $p=3$, $N=n+3k$, and write

\[
W(n,k)=\frac{N!}{k!},
\qquad
H(n,k)=\frac{C_4(n,k)}{W(n,k)}.
\]

The original row is log-concave exactly when

\[
H(n,k)^2\ge T(n,k)H(n,k-1)H(n,k+1),
\tag{4.1}
\]

where

\[
T(n,k)=\frac{W(n,k-1)W(n,k+1)}{W(n,k)^2}
=\frac{k}{k+1}\prod_{i=1}^{3}\frac{N+i}{N-3+i}.
\tag{4.2}
\]

\textbf{Theorem 4.1.} For every $n\ge0$, the row $C_4(n,\cdot)$ is log-concave.

\textbf{Proof.} Assume the preceding row $x_j=H(n-1,j)$ satisfies (4.1). First work on four
consecutive entries in its positive support and set

\[
\rho_j=\frac{x_j}{x_{j-1}}.
\]

At positive entries, two adjacent weighted inequalities give

\[
\rho_{k-1}\ge T(n-1,k-1)\rho_k,
\qquad
\rho_{k+1}\le\frac{\rho_k}{T(n-1,k)}.
\tag{4.3}
\]

The recurrence (2.2) has

\[
a_j=\frac{n+3j-1}{n+3j},
\qquad
b_j=\frac{j}{n+3j}.
\]

After dividing the desired defect in row $n$ by $x_{k-1}^2$, its right-hand side is
increasing in $1/\rho_{k-1}$ and in $\rho_{k+1}$. Thus the two equality cases in (4.3)
give the worst possible value. With $\rho=\rho_k$, it is enough to prove

\[
(a_k\rho+b_k)^2\ge T(n,k)\rho
\left(a_{k-1}+\frac{b_{k-1}}{T(n-1,k-1)\rho}\right)
\left(\frac{a_{k+1}\rho}{T(n-1,k)}+b_{k+1}\right).
\tag{4.4}
\]

Put $M=n+3k$. Clearing the positive denominator

\[
M^2(M-3)^2(M-2)^2(M-1)^2
\]

turns the left side minus the right side of (4.4) into a quadratic polynomial
$P(\rho)$. Its leading coefficient is

\[
9(M-3)^2(M-2)^2(M-1)^2>0,
\]

and its discriminant is

\[
-972k^2(M-4)^2(M-3)^2(M-2)^3(M-1)^2(M+2)<0.
\tag{4.5}
\]

Hence $P(\rho)>0$ for every $\rho>0$. This proves the internal inequalities for
$n\ge3$ and $2\le k\le n-1$. The case $k=1$ is immediate because $C_4(n,0)=0$.
At the left endpoint, the same expanded inequality is obtained as the limit
$x_0\downarrow0$, for which $1/\rho_1\downarrow0$; this remains below the upper bound in
(4.3). At the right endpoint, the limit $x_n\downarrow0$ gives $\rho_n\downarrow0$, again
no larger than the worst-case substitution. The short
initial rows are direct. Induction proves the theorem. $\square$

The negative discriminant in (4.5) is specific to the weighted cone at $p=3$. At $p=4$
the analogous discriminant is positive, and actual rows enter the interval where the
quadratic is negative. The fifth-order case therefore needs a different argument.

\section{The sharp cutoff after order five}


The first nontrivial three-term row already determines the failure side of the classification.

\textbf{Proposition 5.1.} For $r\ge2$, the row $C_r(3,\cdot)$ is log-concave if and only if
$2\le r\le5$. In particular, the assertion that every row is log-concave fails for every
$r\ge6$.

\textbf{Proof.} The three positive entries are

\[
C_r(3,1)=(r+1)!,
\]

\[
C_r(3,2)=\frac{(2r+1)!}{r(r+1)},
\]

and

\[
C_r(3,3)=\frac{(3r)!}{6r^3}.
\]

Thus the only nontrivial log-concavity ratio is

\[
R(r)=\frac{C_r(3,2)^2}{C_r(3,1)C_r(3,3)}
=\frac{6r(2r+1)!^2}{(r+1)^2(3r)!(r+1)!}.
\tag{5.1}
\]

Its adjacent ratio is

\[
\frac{R(r+1)}{R(r)}
=\frac{4(r+1)^4(2r+3)^2}{3r(r+2)^3(3r+1)(3r+2)}.
\tag{5.2}
\]

The denominator of (5.2) minus its numerator equals

\[
11r^6+77r^5+168r^4+80r^3-136r^2-144r-36.
\]

After substituting $r=s+1$, this becomes

\[
11s^6+143s^5+718s^4+1742s^3+2047s^2+947s+20,
\]

which is positive for $s\ge0$. Hence $R(r)$ is strictly decreasing for $r\ge1$. Finally,

\[
R(5)=\frac{55}{39}>1,
\qquad
R(6)=\frac{5148}{5831}<1.
\]

This proves the proposition. $\square$

Deb and Sokal's positive conjecture stops at $r=5$. Proposition 5.1 does not refute a
claim they made for $r\ge6$; it shows that their stated endpoint is the sharp endpoint of
the all-row phenomenon.

\section{Reduction of the fifth-order case}


It remains to prove the last positive layer.

\textbf{Theorem 6.1.} For every $n\ge0$, the row $C_5(n,\cdot)$ is log-concave. At every
internal index for which the three terms are positive, the inequality is strict.

Write

\[
d=n-k,
\qquad
f(z)=\sum_{j\ge0}\frac{z^j}{j+5},
\qquad
B(k,d)=[z^d]f(z)^k.
\tag{6.1}
\]

The exponential formula for permutations by cycles gives

\[
C_5(n,k)=\frac{(5k+d)!}{k!}B(k,d).
\tag{6.2}
\]

For $k\ge2$, put

\[
\mu=\frac{d}{k},\qquad q=\frac1k.
\]

After substituting (6.2) at $k-1,k,k+1$, the required inequality is equivalent to

\[
B(k,d)^2\ge T_k(\mu)B(k-1,d+1)B(k+1,d-1),
\tag{6.3}
\]

where

\[
T_k(\mu)=\frac{1}{1+q}
\prod_{i=1}^{4}
\frac{\mu+5+iq}{\mu+5-(i-1)q}.
\tag{6.4}
\]

The proof of (6.3) divides the parameter plane into two exact wedges and one compact band.

\section{Two exact wedges}


Define

\[
Q_d(k)=d!\,5^kB(k,d).
\tag{7.1}
\]

The identity

\[
(1-z)\bigl(zf'(z)+5f(z)\bigr)=1
\]

implies

\[
Q_d(k)=
\frac{d(5k+d-1)}{5k+d}Q_{d-1}(k)
+\frac{5k}{5k+d}Q_d(k-1).
\tag{7.2}
\]

To state the two-step lemma with the correct row index, put

\[
q_n(k)=Q_{n-k}(k).
\tag{7.3}
\]

Then (7.2) is the triangular recurrence

\[
q_n(k)=A(n,k)q_{n-1}(k)+D(n,k)q_{n-1}(k-1),
\tag{7.4}
\]

where

\[
A(n,k)=\frac{(n-k)(n+4k-1)}{n+4k},
\qquad
D(n,k)=\frac{5k}{n+4k}.
\tag{7.5}
\]

\textbf{Lemma 7.1 (two-step preservation).} Two successive applications of (7.4), from the
row indexed by $n-2$ to the row indexed by $n$, preserve
log-concavity at every internal index with $k\ge2$ and $d=n-k\ge3$. The boundary
defects corresponding to $d=1,2$ are strictly positive. Starting from the short initial
rows, every row $q_n(\cdot)$ is therefore log-concave.

\textbf{Proof certificate.} Write the two-step map as

\[
y_k=\alpha_kx_k+\beta_kx_{k-1}+\gamma_kx_{k-2}.
\]

For a log-concave input, write four consecutive ratios as $u\ge v\ge w\ge t$. Direct
differentiation of the output defect shows that its minimum occurs at $u=v$ and $t=w$.
Put $v=w+h$ with $h\ge0$. After substituting $k=K+2$ and $d=D+3$, the numerator of the
defect has exactly twelve monomials in $(h,w)$. Eleven coefficients are coefficientwise
nonnegative in $(K,D)$. In the $h$-free part, the only coefficient not individually
nonnegative is the coefficient $c_1$ of $w$; the exact identity

\[
4c_0c_2-c_1^2\ge0
\]

is coefficientwise nonnegative in $(K,D)$. Since the coefficient $c_2$ and the common
denominator are positive, the quadratic $c_0+c_1w+c_2w^2$ is nonnegative for $w>0$.
This establishes the internal part of the lemma over exact rational functions. For
$d=1,2$, the defects are respectively

\[
\frac{5(33k+2)}{252}
\]

and

\[
\frac{25(1155k^3+2426k^2+1435k+168)}{63504}.
\]

The full twelve-coefficient expansion and the shifted polynomial
$4c_0c_2-c_1^2$ are regenerated exactly by the verification program listed in
Section 11. This is a finite symbolic certificate for the displayed algebraic reduction,
not a sample of $(n,k)$ values.

We next connect Lemma 7.1 to (6.3). Since $d=n-k$,

\[
q_n(k)^2\ge q_n(k-1)q_n(k+1)
\]

is exactly

\[
Q_d(k)^2\ge Q_{d+1}(k-1)Q_{d-1}(k+1).
\tag{7.6}
\]

Using (7.1), inequality (6.3) follows from (7.6) whenever

\[
W(k,d)=\frac{d}{d+1}T_k(d/k)\le1.
\tag{7.7}
\]

Explicitly,

\[
W(k,d)=
\frac{dk(d+5k+1)(d+5k+2)(d+5k+3)(d+5k+4)}
{(d+1)(k+1)(d+5k)(d+5k-1)(d+5k-2)(d+5k-3)}.
\tag{7.8}
\]

Let $P(k,d)$ be the numerator of $1-W(k,d)$ after clearing the positive denominator.
The three substitutions needed below give polynomials with nonnegative coefficients:

\[
\begin{aligned}
P(2(U+1)+1+S,U+1)
={}&625S^5+5375S^4U+8375S^4\\
&+17450S^3U^2+54100S^3U+41925S^3\\
&+25630S^2U^3+117810S^2U^2+180885S^2U+93025S^2\\
&+15125SU^4+89540SU^3+200255SU^2+201630SU+77450S\\
&+1331U^5+6655U^4+11495U^3+7865U^2+2054U+600,
\end{aligned}
\tag{7.9}
\]

\[
\begin{aligned}
P(2(U+12),U+12)
={}&1331U^5+64735U^4+1185415U^3\\
&+9744665U^2+31222854U+8690040,
\end{aligned}
\tag{7.10}
\]

and, for $k=K+3$, $d=12(K+3)+S$,

\[
\begin{aligned}
P(k,d)
={}&142477K^5+93925K^4S+1754230K^4
+18530K^3S^2\\
&+1032920K^3S+8228255K^3+1610K^2S^3
+158220K^2S^2\\
&+4225545K^2S+17796230K^2+65KS^4+9320KS^3\\
&+449165KS^2+7608500KS+16368048K\\
&+S^5+190S^4+13475S^3+423830S^2+5076504S+3630240.
\end{aligned}
\tag{7.11}
\]

Thus $W<1$ for $k\ge2d+1$, for the boundary $k=2d$ when $d\ge12$, and for
$d\ge12k$ when $k\ge3$. The remaining points on $k=2d$ have $d\le11$ and are covered
by exact all-$k$ strip certificates through $d=18$.

The low-density wedge also needs the internal index $k=2$, which is not supplied by
(7.11). We record its direct proof.

\textbf{Lemma 7.2 (the $k=2$ line).} The fifth-order row inequality at $k=2$ holds for every
$n\ge3$.

\textbf{Proof.} Put

\[
A=\sum_{j=5}^{n+3}\frac1j,
\qquad
e=\sum_{j=n+4}^{n+7}\frac1j.
\]

Cycle decomposition gives

\[
C_5(n,1)=(n+3)!,
\qquad
C_5(n,2)=(n+7)!\,A,
\]

and

\[
C_5(n,3)=(n+11)!\,S_3,
\]

where

\[
S_3=\sum_{\substack{5\le i<j\le n+7\\j-i\ge5}}\frac1{ij}.
\]

Dropping the separation condition gives

\[
S_3<\frac{(A+e)^2}{2}.
\]

Consequently, the desired inequality follows from

\[
2A^2\ge R(A+e)^2,
\qquad
R=\prod_{j=0}^{3}\frac{n+8+j}{n+4+j}.
\tag{7.12}
\]

Set $x=n+4$. For $x\ge29$,

\[
A\ge\log(x/5),
\qquad
e\le\frac4x,
\qquad
\sqrt R\le\left(1+\frac4x\right)^2.
\]

The two factors in

\[
\left(1+\frac4x\right)^2
\left(1+\frac{4}{x\log(x/5)}\right)
\]

decrease with $x$. At $x=29$, use $\log(29/5)>7/4$ and $\sqrt2>7/5$ to obtain

\[
\left(\frac{33}{29}\right)^2
\left(1+\frac{16}{203}\right)
=\frac{238491}{170723}<\frac75<\sqrt2.
\]

This is (7.12) for every $n\ge25$. The logarithmic bound follows from a rational Taylor
upper bound for $\exp(7/4)$; the cases $3\le n\le24$ are checked exactly over the
rationals. $\square$

\textbf{Lemma 7.3 (exact wedges).} Inequality (6.3) holds strictly in each of the following
regions:

\[
k\ge2d,
\tag{7.13}
\]

and

\[
d\ge12k.
\tag{7.14}
\]

For the first wedge, (7.9) handles $k\ge2d+1$, (7.10) handles $k=2d$ for $d\ge12$,
and the strip certificates handle the remaining boundary points. For the second wedge,
(7.11) handles $k\ge3$, $k=1$ is a zero-boundary inequality, and Lemma 7.2 handles
$k=2$. In each symbolic region $W<1$, so positivity of the neighboring $Q$-entries makes
the original inequality strict. All boundary exceptions are therefore included rather than
inferred from an asymptotic estimate.

The only region left after Lemma 7.3 is

\[
\frac12<\mu<12.
\tag{7.15}
\]

\section{A common saddle point and the compact-band margin}


For each $\mu$ in (7.5), there is a unique $z\in(0,1)$ satisfying

\[
\mu=\frac{zf'(z)}{f(z)}.
\tag{8.1}
\]

Define a span-one probability distribution by

\[
\Pr_z(J=j)=\frac{z^j}{(j+5)f(z)},\qquad j\ge0.
\tag{8.2}
\]

Let $\sigma^2$ and $\kappa_s$ be its variance and cumulants. If $S_m$ is a sum of $m$
independent copies of $J$, then

\[
B(k,d)=f(z)^kz^{-d}\Pr_z(S_k=d).
\tag{8.3}
\]

The same $z$ works for all three coefficients in (6.3). The central term $S_k=d$ is at
its mean. The other two probabilities have displacements

\[
S_{k-1}=d+1:\quad \mu+1,
\qquad
S_{k+1}=d-1:\quad -(\mu+1)
\tag{8.4}
\]

from their respective means. All exponential saddle factors in (8.3) therefore cancel.

Put

\[
\Delta(k,\mu)=
\log\frac{B(k,d)^2}
{T_k(\mu)B(k-1,d+1)B(k+1,d-1)}.
\tag{8.5}
\]

The first-order coefficient in the common-saddle expansion is

\[
M(\mu)=\frac{(\mu+1)^2}{\sigma^2}+1-\frac{16}{\mu+5}.
\tag{8.6}
\]

It is positive precisely because the saddle variance satisfies

\[
\sigma^2(11-\mu)<(\mu+1)^2(\mu+5).
\tag{8.7}
\]

The verification of (8.7) reduces, after writing $f$ as its elementary logarithmic
expression, to a rational inequality whose differentiated remainder is a nonnegative
square. On the full compact band, a rational Bernstein certificate sharpens it to

\[
M(\mu)>\frac3{100}.
\tag{8.8}
\]

Expanding the three local probabilities through order five gives

\[
\Delta(k,\mu)=
\frac{M}{k}+\frac{C_2}{k^2}+\frac{C_3}{k^3}
+\frac{C_4}{k^4}+\frac{C_5}{k^5}+R_k.
\tag{8.9}
\]

Each $C_j$ is a rational expression in $z$, $f(z)$, and the cumulants through order ten.
Exact two-variable Bernstein certificates on the complete saddle band prove

\[
C_2>-2M,
\qquad
C_3>0,
\qquad
C_4>-4M,
\qquad
C_5>0.
\tag{8.10}
\]

Hence the retained part of (8.9) is at least

\[
\frac3{100}\left(\frac1k-\frac2{k^2}-\frac4{k^4}\right).
\tag{8.11}
\]

\section{Effective Fourier control for the infinite tail}


We now give the effective local expansion used to bound $R_k$. Let

\[
\phi_z(t)=\frac{f(ze^{it})}{f(z)},
\qquad
\psi_z(t)=e^{-i\mu t}\phi_z(t),
\]

and normalize the lattice probabilities by

\[
\overline p_m(b)=
\sigma\sqrt{2\pi m}\,\Pr_z(S_m=m\mu+b).
\tag{9.1}
\]

Fourier inversion with

\[
u=\sigma\sqrt m\,t,
\qquad
\varepsilon=m^{-1/2}
\]

gives

\[
\overline p_m(b)=\frac1{\sqrt{2\pi}}
\int_{-\pi\sigma\sqrt m}^{\pi\sigma\sqrt m}
e^{-ibu/(\sigma\sqrt m)}
\psi_z\!\left(\frac{u}{\sigma\sqrt m}\right)^m\,du.
\tag{9.2}
\]

The Gaussian term must be extracted before estimating the correction. On the local
analytic disk,

\[
e^{-ibu/(\sigma\sqrt m)}
\psi_z\!\left(\frac{u}{\sigma\sqrt m}\right)^m
=e^{-u^2/2}
\exp\left(\sum_{j\ge1}E_j(u,b)\varepsilon^j\right),
\tag{9.3}
\]

where

\[
E_1(u,b)=-\frac{ibu}{\sigma}+\frac{\kappa_3(iu)^3}{6\sigma^3},
\]

and, for $2\le j\le8$,

\[
E_j(u,b)=\frac{\kappa_{j+2}(iu)^{j+2}}{(j+2)!\sigma^{j+2}}.
\tag{9.4}
\]

The three displacements needed in (8.5) satisfy $|b|\le\mu+1$. Exact rational
certificates over the saddle band prove

\[
\frac{\mu+1}{\sigma}<\frac74
\]

and the termwise bounds

\[
|E_j(u,b)|\le A_j(u),
\tag{9.5}
\]

where

\[
\begin{aligned}
A_1&=\frac74|u|+\frac{11}{20}|u|^3,&
A_2&=\frac{13}{20}|u|^4,\\
A_3&=\frac{17}{20}|u|^5,&
A_4&=\frac65|u|^6,\\
A_5&=\frac74|u|^7,&
A_6&=\frac52|u|^8,\\
A_7&=\frac{15}{4}|u|^9,&
A_8&=\frac{11}{2}|u|^{10}.
\end{aligned}
\tag{9.6}
\]

Define the exact local polynomials $H_s(u,b)$ by

\[
\exp\left(\sum_{j=1}^{8}E_j(u,b)\varepsilon^j\right)
=\sum_{s\ge0}H_s(u,b)\varepsilon^s.
\]

Equating coefficients gives the recurrence

\[
sH_s=\sum_{j=1}^{\min(8,s)}jE_jH_{s-j}.
\tag{9.7}
\]

The five probability coefficients used in Section 8 are therefore

\[
p_j(b)=\frac1{\sqrt{2\pi}}
\int_{\mathbb R}e^{-u^2/2}H_{2j}(u,b)\,du,
\qquad 1\le j\le5.
\tag{9.8}
\]

The odd integrated terms vanish. Substituting $b=0,\mu+1,-(\mu+1)$ at sample sizes
$k,k-1,k+1$, expanding $(k\pm1)^{-j}$ in $1/k$, and including the normalizing and
target-weight terms produces exactly the coefficients $M,C_2,\ldots,C_5$ in (8.9).

To control cumulants above order ten, write

\[
5f(z)=\frac1{1-C(z)},
\qquad
C(z)=\sum_{\ell\ge1}c_\ell z^\ell,
\]

where $c_\ell>0$ and $\sum_{\ell\ge1}c_\ell=1$. Define

\[
\log(5f(z))=\sum_{\ell\ge1}\beta_\ell z^\ell,
\]

For completeness, both assertions about $C$ follow directly from the defining function.
Differentiation gives

\[
z(1-z)C'(z)=5(1-C(z))(z-C(z)).
\]

Writing $C(z)=\sum_{m\ge1}c_mz^m$ yields

\[
(m+5)c_m=(m-6)c_{m-1}
+5\sum_{i=1}^{m-1}c_ic_{m-i},
\qquad m\ge2.
\tag{9.9}
\]

The first five coefficients are

\[
\frac56,\quad \frac5{252},\quad \frac5{378},\quad
\frac{605}{63504},\quad \frac{1381}{190512}.
\]

They are positive. For $m=6$ the convolution in (9.9) is positive, and for $m\ge7$
both terms on the right are nonnegative with a positive convolution. Induction proves
$c_m>0$ for every $m$. Finally, $5f(z)\to\infty$ as $z\uparrow1$, so
$C(z)=1-1/(5f(z))\uparrow1$; hence $\sum_{m\ge1}c_m=1$.

Moreover,

\[
z\frac{d}{dz}\log(5f(z))=\frac{5(z-C(z))}{1-z}.
\]

Comparing coefficients gives

\[
\ell\beta_\ell=5\sum_{j>\ell}c_j,
\qquad
0<\beta_\ell\le\frac5\ell.
\tag{9.10}
\]

Thus (8.2) is compound Poisson. Together with $(1-z)\sigma>8/15$, equation (9.10)
gives the all-order envelope as follows. The compound-Poisson cumulants satisfy

\[
\kappa_s=\sum_{\ell\ge1}\ell^s\beta_\ell z^\ell
\le5\sum_{\ell\ge1}\ell^{s-1}z^\ell
\le\frac{5(s-1)!}{(1-z)^s}.
\]

After division by $s!\sigma^s$, this is

\[
0<\frac{\kappa_s}{s!\sigma^s}
\le\frac5s\left(\frac{15}{8}\right)^s.
\tag{9.11}
\]

We also need two characteristic-function estimates. The integral representation

\[
f(w)=\int_0^1\frac{x^4}{1-wx}\,dx
\]

implies that $|\phi_z(t)|$ decreases strictly on $[0,\pi]$. Indeed, after symmetrizing
the double integral and putting $q=zx$, $r=zy$, and $y_0=\cos t$, the derivative of
the kernel with respect to $y_0$ has numerator

\[
(1-q)(1-r)\mathcal Q(q,r)
+8qr(1-q)(1-r)(1-y_0)+4qr(q+r)(1-y_0)^2,
\]

where

\[
\mathcal Q(q,r)=
(q+r)\left(1-\frac{q+r}{2}\right)^2
+\frac{(q-r)^2(8-q-r)}4\ge0.
\]

The same rational saddle-band certificate proves, for the centered fourth moment
$\mu_4$,

\[
\mu_4<18\sigma^4.
\tag{9.12}
\]

For an independent copy $J'$, this gives

\[
\mathbb E(J-J')^4=2\mu_4+6\sigma^4<42\sigma^4.
\]

Using $1-\cos x\ge x^2/2-x^4/24$, for $|t|\le1/(2\sigma)$ we obtain

\[
1-|\phi_z(t)|^2
\ge\sigma^2t^2-\frac{42\sigma^4t^4}{24}
\ge\frac9{16}\sigma^2t^2.
\tag{9.13}
\]

Hence, in the variable $u=\sigma\sqrt m\,t$,

\[
|\phi_z(t)|^m\le e^{-9u^2/32}.
\tag{9.14}
\]

At $t_0=1/(2\sigma)$, (9.13) gives

\[
1-|\phi_z(t_0)|^2\ge\frac9{64}.
\]

Monotonicity therefore yields

\[
|\phi_z(t)|^2\le\frac{55}{64},
\qquad
t_0\le|t|\le\pi.
\tag{9.15}
\]

Now take the moving cutoff

\[
U(m)=6\left(\frac{m}{1000}\right)^{1/16}.
\tag{9.16}
\]

For a polynomial $P$, write

\[
\langle P\rangle_G=
\frac1{\sqrt{2\pi}}\int_{\mathbb R}e^{-u^2/2}|P(u)|\,du.
\]

Let $\widehat H_s(u)$ be the positive formal majorants generated by

\[
\exp\left(\sum_{j=1}^{8}A_j(u)\varepsilon^j\right)
=\sum_{s\ge0}\widehat H_s(u)\varepsilon^s.
\]

They satisfy the same recurrence (9.7) with $E_j$ replaced by $A_j$, and dominate the
absolute coefficients of the exact $H_s(u,b)$.

The normalized error in (9.2) is divided into the following six explicit majorants:

\[
E_{\mathrm{core}}(m)=
\sum_{\substack{12\le s\le40\\s\ \mathrm{even}}}
m^{-s/2}\langle \widehat H_s\rangle_G,
\tag{9.17}
\]

\[
E_{\mathrm{point}}(m)\le
\frac{4U(m)}5
\left[
\exp\left(\sum_{j=1}^{8}A_j(U(m))m^{-j/2}\right)
-\sum_{s=0}^{40}\widehat H_s(U(m))m^{-s/2}
\right],
\tag{9.18}
\]

and

\[
E_{\ge11}(m)\le
\frac{20}{11(1-q_m)}
\left(\frac{15}{8}\right)^{11}
m^{-9/2}\mathbb E|Z|^{11},
\qquad
q_m=\frac{15U(m)}{8\sqrt m}.
\tag{9.19}
\]

Here $Z$ is standard normal. The retained polynomial outside the core is bounded by

\[
E_{\mathrm{poly}}(m)\le
\sum_{j=0}^{5}m^{-j}
\sum_{\ell}|h_{2j,\ell}|
\frac{2}{\sqrt{2\pi}}\int_{U(m)}^\infty
u^\ell e^{-u^2/2}\,du,
\tag{9.20}
\]

where $\widehat H_{2j}(u)=\sum_\ell h_{2j,\ell}u^\ell$. On $6\le|u|\le8$, the certified
cumulant bounds sharpen (9.14) to $e^{-111u^2/250}$. Thus the remaining local Fourier
tail is at most

\[
E_{\mathrm{near}}(m)\le
\frac{2}{\sqrt{2\pi}}
\left[
\int_{U(m)}^{8}e^{-111u^2/250}\,du
+\int_{\max(U(m),8)}^\infty e^{-9u^2/32}\,du
\right],
\tag{9.21}
\]

with an empty first integral when $U(m)\ge8$. Finally, (9.15) and $\sigma<17$ give

\[
E_{\mathrm{far}}(m)\le
\sigma\sqrt{2\pi m}\left(\frac{55}{64}\right)^{m/2}.
\tag{9.22}
\]

All exponentials and Gaussian tails in (9.17)--(9.22) are bounded by rational Taylor or
Mills majorants in the certificate. At $m=1000$, the contributions to a bound of the
form $C/m$ are respectively

\begin{center}
\renewcommand{\arraystretch}{1.08}
\begin{tabular}{@{}l r@{}}
\toprule
component & contribution to $C$ \\
\midrule
$E_{\mathrm{core}}$ & $0.0003482608957161$ \\
$E_{\mathrm{point}}$ & $0.0005681456909691$ \\
$E_{\ge11}$ & $0.0002847773304812$ \\
$E_{\mathrm{poly}}$ & $0.0002761906358576$ \\
first integral in $E_{\mathrm{near}}$ & $0.0000171697025750$ \\
second integral in $E_{\mathrm{near}}$ & $0.0002701459962718$ \\
$E_{\mathrm{far}}$ & $<2.1\times10^{-27}$ \\
\bottomrule
\end{tabular}
\end{center}

The displayed entries are decimal approximations to exact rational upper bounds; their
exact sum is $0.001764690251870703\ldots<0.001764690252$.

The moving cutoff makes these endpoint bounds uniform. For (9.17), multiplying by $m$
leaves powers $m^{1-s/2}$ with $s\ge12$. For the point tail, every omitted order
$s\ge41$ has degree at most $3s$, so its worst scaling exponent after multiplication by
$m$ is

\[
\frac{17}{16}-\frac{5s}{16}<0.
\]

In (9.19), $q_m$ is proportional to $m^{-7/16}$. For the polynomial Gaussian tail,
the monomial $m^{-j}u^\ell$ has $\ell\le6j$. After multiplication by $m$, its worst
logarithmic scaling exponent at the base cutoff $U=6$ is

\[
1-j+\frac{\ell-1}{16}-\frac{36}{16}
\le\frac{-21-10j}{16}<0.
\]

The Mills bounds for (9.21) have the form $mU^{-1}e^{-cU^2}$ and decrease from
$U=6$ or $U=8$; the transition occurs at

\[
m=\left\lceil1000(4/3)^{16}\right\rceil.
\]

The far bound decreases because its exponential factor dominates the square-root
prefactor. These exact exponent checks prove

\[
\left|
\overline p_m(b)-\left(1+\sum_{j=1}^{5}\frac{p_j(b)}{m^j}\right)
\right|
\le\frac{C}{m},
\qquad
C<0.001764690252,
\tag{9.23}
\]

for every $m\ge1000$ and each of the three displacements.

The retained probability polynomial differs from one by less than $0.009085849$, so the
passage to logarithms increases the absolute probability error by less than a factor two.
The absolute weights in the logarithmic combination are $2+1+1=4$, and the smallest
sample size is $k-1$; the Fourier contribution is therefore at most $8C/(k-1)$. The
tails from the logarithm series, re-expansion of $(k\pm1)^{-j}$, normalizing factors, and
(6.4) all begin at order $k^{-6}$ and are separately bounded by rational geometric
majorants. At $k=1001$ their total, together with the Fourier contribution, is

\[
|R_k|<1.419319083\times10^{-5}.
\tag{9.24}
\]

The right side of (8.11) is greater than

\[
2.991014967\times10^{-5}.
\tag{9.25}
\]

After multiplication by $k$, every error term just listed is nonincreasing, whereas
(8.11) is increasing. Therefore $\Delta(k,\mu)>0$ throughout the compact band for every
$k\ge1001$.

\section{Certified finite prefix and completion of the proof}


It remains to cover the compact band for $2\le k\le1000$. Let

\[
c(k,d)=5^kB(k,d).
\]

It satisfies the positive recurrence

\[
(d+5k)c(k,d)=(d-1+5k)c(k,d-1)+5k\,c(k-1,d),
\qquad
c(k,0)=1.
\tag{10.1}
\]

A 40-digit decimal computation with an explicit relative-error recurrence checks every
point

\[
2\le k\le1000,
\qquad
\frac{k}{2}<d<12k.
\tag{10.2}
\]

There are $5,754,989$ such inequalities. The propagated relative distortion is less than
$10^{-30}$. The smallest computed relative margin is

\[
0.000032205417799388957039776739958609472
\]

at $(k,d)=(1000,1868)$, which is larger than both the rounding budget and the acceptance
threshold $10^{-25}$. Thus every inequality in (10.2) is strict.

We can now prove Theorem 6.1. The region $k\ge2d$ is covered by (7.13), and the region
$d\ge12k$ by (7.14). The remaining compact band is covered by the certified finite prefix
when $k\le1000$ and by the effective Fourier estimate when $k\ge1001$. The cases
$k=0,1$, $d=0$, and the right boundary are elementary; $k=2$ is included by its
independent exact proof. These regions contain every admissible pair $(n,k)$ and give
strictness whenever all three neighboring terms are positive. $\square$

Combining Theorems 3.1, 4.1, and 6.1 with the known $r=1,2$ cases proves the positive
direction of Theorem 1.1. Proposition 5.1 proves the converse.

\section{Reproducibility and proof boundary}


The proof has three kinds of evidence, which should not be conflated.

\begin{enumerate}
\item Sections 2--5 are ordinary symbolic proofs. The accompanying scripts independently
regenerate the rational identities and coefficient signs.
\item Lemmas 7.1--7.3 and the coefficient inequalities (8.8)--(8.10) are finite symbolic
certificates. Their programs reconstruct the relevant rational functions and verify all
shifted monomial or Bernstein coefficients exactly.
\item Section 10 is an exhaustive finite computation with a proved rounding budget. It is
used only for the explicit prefix (10.2). Section 9 proves the infinite remainder.

\end{enumerate}

The complete internal verification map, interpreter versions, commands, expected terminal
markers, and SHA-256 hashes are recorded in
\texttt{verification/README.md}. The entry points are:

\begin{itemize}
\item \texttt{verification/\allowbreak{}verify\_\allowbreak{}normalized\_\allowbreak{}logconcavity.py};
\item \texttt{verification/\allowbreak{}verify\_\allowbreak{}r4\_\allowbreak{}weighted\_\allowbreak{}cone.py};
\item \texttt{verification/\allowbreak{}verify\_\allowbreak{}two\_\allowbreak{}step\_\allowbreak{}lc\_\allowbreak{}preserver.py};
\item \texttt{verification/\allowbreak{}verify\_\allowbreak{}small\_\allowbreak{}excess\_\allowbreak{}strips.py};
\item \texttt{verification/\allowbreak{}verify\_\allowbreak{}sparse\_\allowbreak{}cycle\_\allowbreak{}k2.py};
\item \texttt{verification/\allowbreak{}verify\_\allowbreak{}saddle\_\allowbreak{}curvature.py};
\item \texttt{verification/\allowbreak{}verify\_\allowbreak{}bulk\_\allowbreak{}edgeworth\_\allowbreak{}margin.py};
\item \texttt{verification/\allowbreak{}verify\_\allowbreak{}bulk\_\allowbreak{}second\_\allowbreak{}order\_\allowbreak{}margin.py};
\item \texttt{verification/\allowbreak{}verify\_\allowbreak{}bulk\_\allowbreak{}third\_\allowbreak{}order\_\allowbreak{}margin.py};
\item \texttt{verification/\allowbreak{}verify\_\allowbreak{}bulk\_\allowbreak{}fourth\_\allowbreak{}order\_\allowbreak{}margin.py};
\item \texttt{verification/\allowbreak{}verify\_\allowbreak{}bulk\_\allowbreak{}fifth\_\allowbreak{}order\_\allowbreak{}margin.py};
\item \texttt{verification/\allowbreak{}verify\_\allowbreak{}renewal\_\allowbreak{}reformulation.py};
\item \texttt{verification/\allowbreak{}verify\_\allowbreak{}low\_\allowbreak{}cumulant\_\allowbreak{}envelope.py};
\item \texttt{verification/\allowbreak{}verify\_\allowbreak{}characteristic\_\allowbreak{}tail\_\allowbreak{}constants.py};
\item \texttt{verification/\allowbreak{}verify\_\allowbreak{}effective\_\allowbreak{}fourier\_\allowbreak{}remainder.py};
\item \texttt{verification/\allowbreak{}verify\_\allowbreak{}bulk\_\allowbreak{}finite\_\allowbreak{}prefix.py}.

\end{itemize}

The reference environment is CPython 3.13.11 with SymPy 1.14.0. The aggregate PowerShell
runner accepts an explicit Python interpreter and executes every program from the published
\texttt{verification} directory.
The fifth-order Bernstein certificate is the slowest symbolic step and required about
27 minutes in the recorded audit environment. Runtime is not part of the mathematical
claim.

The machine checks establish the stated algebraic and finite inequalities, conditional on
the correctness of the reductions and implementations. Sections 7 and 9 state the human
reductions that connect those inequalities to Theorem 6.1; the programs reconstruct their
rational certificates. These remain separate audit obligations. No claim of real-rootedness,
total positivity, or Hankel-total positivity for the $r=3,4,5$ rows is made here.

\section{AI-assisted research and writing disclosure}


Carptopus is the author and is responsible for the mathematical claims, source selection,
verification materials, and publication decisions. OpenAI Codex was used to assist with
symbolic exploration, exact verification programs, proof checking, literature searches, and
drafting. All theorem statements and proof dependencies were reviewed against the archived
source and verification materials before this draft was prepared.

\begin{thebibliography}{9}

\bibitem{deb-sokal}
B.~Deb and A.~D. Sokal,
\emph{Higher-order Stirling cycle and subset triangles: Total positivity,
continued fractions and real-rootedness},
arXiv:2507.18959v1 (2025).
\href{https://arxiv.org/abs/2507.18959}{arXiv:2507.18959}.

\bibitem{bcc30}
P.~J. Cameron, ed.,
\emph{Problems from BCC30}, Problem 30.9,
arXiv:2409.07216 (2024).
\href{https://arxiv.org/abs/2409.07216}{arXiv:2409.07216}.

\bibitem{sagan}
B.~E. Sagan,
\emph{Inductive and injective proofs of log concavity results},
Discrete Mathematics \textbf{68} (1988), 281--292.
\href{https://doi.org/10.1016/0012-365X(88)90116-4}{doi:10.1016/0012-365X(88)90116-4}.

\end{thebibliography}

\end{document}
